\centerline{\bf \Huge  Матдрака}
\title{Матдрака}

\begin{flushright}
    \scriptsize Лучший друг не пытается разнять вашу драку\\
    Он влетает в драку с ударом ноги.\\
    Джейсон Стэтхэм.
\end{flushright}

1. Маша пробежала 1 км со средней скоростью 4 м/с. С какой средней скоростью пробежал эту дистанцию Вася, если, стартовав на 25 секунд позже Маши, он финишировал на 25 секунд раньше?

2. Разделите число 80 на две части так, чтобы одна часть составляла 60\% от другой части.

3. Сколько существует трёхзначных чисел, у которых сумма цифр больше произведения цифр?

4. Четыре мецената пожертвовали театру 132 тысячи рублей. При этом второй пожертвовал вдвое больше первого, третий — втрое больше второго, четвёртый — вчетверо больше третьего. Сколько пожертвовал четвёртый?

5. Незнайка лжёт по понедельникам, вторникам и пятницам, а в остальные дни недели говорит правду. В какие дни недели Незнайка может сказать: «Я лгал позавчера и буду лгать послезавтра»?

6. На шахматной доске стоят ладьи так, что каждая из них бьёт N ладей. При каких целых N это возможно? (Ладья бьёт в каждом направлении только ближайшую ладью.)

7. Маша выписывает последовательно на доску по возрастанию все числа, в которых число четных цифр равно числу нечетных цифр. Какое число выпишет Маша 46–м?

8. Три поросенка хранят в жестяной банке красные, желтые и зеленые леденцы. Какое наименьшее число леденцов надо взять наугад из банки так, чтобы каждому поросенку можно было дать по 5 леденцов одного цвета? (У разных поросят леденцы могут быть и разными.)

9. Вася задумал целое число. Коля умножил его не то на 5, не то на 6. Женя прибавил к результату Коли не то 5, не то 6. Саша отнял от результата Жени не то 5, не то 6. В итоге получилось 73. Какое число задумал Вася (перечислите все возможные варианты)?

10. Сумма четырёхзначного натурального числа с его суммой цифр равна 2018. Чему равно само число (необходимо найти все возможные варианты)?

\newpage

11. Сколько на шахматной доске имеется всевозможных прямоугольников, состоящих из четырёх клеток?

12. На какое наибольшее количество прямоугольников можно разрезать (без остатка) по линиям сетки клетчатый квадрат 7×7 так, чтобы среди них не оказалось одинаковых?

13. По двум пересекающимся дорогам с равными постоянными скоростями движутся автомобили Ауди и БМВ. Оказалось, что как в 17:00, так и в 18:00 БМВ находился в два раза дальше от перекрестка, чем Ауди. В какое время Ауди мог проехать перекресток? Укажите все возможные варианты.

14. В деревне Большие Топоры живет 100 детей, а в деревне Средние Топоры — 60, между деревнями проложена прямая дорога длиной 6 км. Посередине между Большими Топорами и Средними расположена деревня Малые Топоры, в которой живет 20 детей. В каком месте нужно построить школу, чтобы суммарное расстояние, которые должны будут каждый день преодолевать школьники, было наименьшим? (В ответе укажите расстояние от школы до каждой из деревень.)

15. Каждый из 12 человек — рыцарь, всегда говорящий правду, или всегда лгущий лжец. Один из них сказал: «Число лжецов среди нас делится на 1», второй: «Число лжецов среди нас делится на 2», …, 12–ый: «Число лжецов среди нас делится на 12». Сколько среди них может быть рыцарей? Укажите все возможные варианты.

16. Бусы — это кольцо, на которое нанизаны бусины. Бусы можно поворачивать и переворачивать, они от этого не меняются. Сколько различных видов бус можно составить из 10 одинаковых красных бусин и двух одинаковых синих бусин?

17. В комнате дед, два отца, два сына и два внука (это дед, отцы, сыновья и внуки людей, находящихся в комнате). Сколько людей могло быть в комнате?

18. Найдите количество прямоугольников, составленных из клеток шахматной доски, которые содержат поле C4. (Одна клетка — это тоже прямоугольник.)

19. Найдите 2009–й член последовательности: 1, 2, 1, 2, 3, 2, 1, 2, 3, 4, 3, 2, 1, 2, 3, 4, 5, 4, 3, 2, 1…
  
20. Маша заменила в примере на умножение двузначных чисел цифры буквами: одинаковые – одинаковыми, разные – разными. У неё получилось АБ·ВГ = БББ Каким мог быть исходный пример? Найдите все возможные варианты.